\chapter{Теоретические и технологические основы построения музыкальных рекомендательных систем}
\label{ch:sequential}

Настоящая глава очерчивает теоретическую рамку, в которой строится экспериментальная часть работы. Изложение организовано по линии «от классической коллаборативной фильтрации к трансформерным моделям последовательных рекомендаций, далее~--- к специфике музыкального домена и контентным признакам». Каждый последующий блок мотивирует архитектурные и методологические решения, принятые в главе~\ref{ch:groups} (групповые агрегаторы) и в экспериментальной главе~\ref{ch:experiment}: выбор SASRec с обычной (plain) бинарной кросс-энтропией в качестве персонального скорера, явное внимание к феномену повторного потребления и адаптация протокола оценки под него, а также использование готовых аудиоэмбеддингов YAMBDA в качестве индивидуального контентного сигнала, который в главе~\ref{ch:groups} поднимается на групповой уровень.

\section{Эволюция подходов к моделированию пользовательских предпочтений}
\label{sec:ch1:evolution}

\subsection{Коллаборативная фильтрация и матричная факторизация}
\label{subsec:ch1:cf-mf}

Исторически первыми подходами к задаче рекомендаций стали методы коллаборативной фильтрации~\cite{koren2009matrix}, исходящие из предположения о том, что пользователи со схожими паттернами потребления будут одинаково оценивать новые элементы. Среди model-based методов де-факто стандартом 2010-х годов стала матричная факторизация (SVD, ALS, BPR~\cite{rendle2009bpr}), в которой каждый пользователь и каждый элемент представляются плотным вектором фиксированной размерности~$d$, а предсказание строится как скалярное произведение:
\begin{equation*}
    \hat{r}_{ui} = \mathbf{p}_u^\top \mathbf{q}_i,
\end{equation*}
где $\mathbf{p}_u \in \mathbb{R}^d$ и $\mathbf{q}_i \in \mathbb{R}^d$~--- латентные представления пользователя и элемента. Подход обладает вычислительной эффективностью и хорошей масштабируемостью, однако игнорирует порядок взаимодействий, что критично для доменов с выраженной временно\'{и} динамикой потребления~--- в первую очередь для музыки.

\subsection{Последовательные рекомендации: марковские цепи, GRU4Rec, Caser}
\label{subsec:ch1:seq-pre-transformer}

Признание роли последовательного характера пользовательского поведения сформировало отдельное направление~--- последовательные рекомендации, в которых история взаимодействий $\mathcal{S}_u = (s_1, s_2, \ldots, s_t)$ рассматривается как упорядоченная последовательность, а задачей модели является предсказание следующего элемента~$s_{t+1}$.

Первые подходы основывались на марковских цепях фиксированного порядка. Затем рекуррентные сети, в частности GRU4Rec~\cite{hidasi2016session}, продемонстрировали способность улавливать более длинные зависимости в пользовательских сессиях. Свёрточные модели типа Caser~\cite{tang2018caser} предложили альтернативный взгляд, трактуя последовательность как <<изображение>> и применяя свёртки разных масштабов. Однако принципиальный скачок качества произошёл с переносом механизма self-attention из обработки естественного языка в задачу последовательных рекомендаций.

\section{Трансформерные модели последовательных рекомендаций}
\label{sec:ch1:transformers}

\subsection{SASRec: каузальный механизм самовнимания}
\label{subsec:ch1:sasrec}

Модель Self-Attentive Sequential Recommendation (SASRec), предложенная Kang и McAuley~\cite{sasrec2018}, стала первой работой, успешно применившей self-attention к задаче последовательных рекомендаций. Авторы отмечают, что марковские модели эффективно захватывают локальные паттерны, тогда как рекуррентные сети способны учитывать долгосрочные зависимости, но страдают от затухания градиентов и последовательности обучения; механизм самовнимания снимает оба ограничения за счёт прямого взвешивания всех элементов истории.

Архитектура SASRec построена на стеке блоков самовнимания. Для входной последовательности $\mathcal{S} = (s_1, \ldots, s_n)$ модель формирует матрицу эмбеддингов $\mathbf{E} \in \mathbb{R}^{n \times d}$, к которой добавляются позиционные эмбеддинги $\mathbf{P} \in \mathbb{R}^{n \times d}$:
\begin{equation*}
    \hat{\mathbf{E}} = \mathbf{E} + \mathbf{P}.
\end{equation*}
Каждый блок применяет scaled dot-product attention:
\begin{equation*}
    \mathrm{Attention}(\mathbf{Q}, \mathbf{K}, \mathbf{V}) = \mathrm{softmax}\!\left(\frac{\mathbf{Q}\mathbf{K}^\top}{\sqrt{d}}\right)\mathbf{V},
\end{equation*}
где $\mathbf{Q}$, $\mathbf{K}$, $\mathbf{V}$~--- линейные проекции входного представления. Для обеспечения авторегрессионного характера предсказания на матрицу логитов накладывается каузальная маска, запрещающая позиции~$i$ обращаться к позициям~$j > i$. SASRec превзошла как рекуррентные (GRU4Rec), так и свёрточные (Caser) бейзлайны на нескольких бенчмарках, обучаясь при этом на порядок быстрее за счёт параллелизации; статья заложила фундамент для всего последующего семейства трансформерных рекомендеров.

\subsection{BERT4Rec и проблема воспроизводимости}
\label{subsec:ch1:bert4rec}

Sun и соавторы~\cite{bert4rec2019} предложили BERT4Rec, адаптирующую идею Cloze-обучения BERT к последовательным рекомендациям: вместо однонаправленного авторегрессионного предсказания модель обучается восстанавливать случайно замаскированные элементы последовательности, используя двунаправленный контекст:
\begin{equation*}
    \mathcal{L} = -\sum_{s_m \in \mathcal{S}_{\mathrm{mask}}} \log P(s_m \mid \tilde{\mathcal{S}}),
\end{equation*}
где $\tilde{\mathcal{S}}$~--- последовательность с заменой части позиций на специальный токен \texttt{[MASK]}. На момент публикации BERT4Rec продемонстрировала выигрыш порядка 7--24\% NDCG@10 над SASRec на четырёх бенчмарках и стала одним из наиболее цитируемых бейзлайнов области.

Последующая работа Petrov и Macdonald~\cite{petrov2022bert4rec_replicability} ставит этот результат под сомнение: при корректной настройке гиперпараметров SASRec оказывается конкурентоспособной с BERT4Rec и даже превосходит её на ряде датасетов. Это указывает на то, что преимущества двунаправленного внимания не безусловны и в значительной мере зависят от характеристик данных и протокола оценки. Содержательным следствием является вывод: на этапе планирования эксперимента имеет смысл предпочесть SASRec в качестве более компактного и воспроизводимого бейзлайна, а наблюдаемые расхождения в качестве искать не на уровне архитектурного выбора attention-блока, а на уровне функции потерь и стратегии негативного сэмплирования. Именно эту линию рассуждения и развивает работа о gSASRec, рассмотренная далее.

\subsection{gSASRec и обобщённая бинарная кросс-энтропия}
\label{subsec:ch1:gsasrec}

Petrov и Macdonald~\cite{gsasrec2023} в работе, получившей награду Best Paper на RecSys 2023, доводят эту линию до архитектурно нейтрального утверждения: значительная часть различий в производительности между SASRec и BERT4Rec обусловлена не блоком внимания, а выбором функции потерь и количеством негативных примеров на один положительный. Авторы предлагают обобщённую бинарную кросс-энтропию (gBCE), параметризованную числом негативных примеров~$t$:
\begin{equation*}
    \mathcal{L}_{\mathrm{gBCE}} = -\log \sigma(\hat{y}_{\mathrm{pos}}) - t \cdot \log\!\left(1 - \sigma(\hat{y}_{\mathrm{neg}})\right),
\end{equation*}
где $\hat{y}_{\mathrm{pos}}$ и $\hat{y}_{\mathrm{neg}}$~--- предсказанные оценки для положительного и негативного примеров, а $\sigma$~--- сигмоида. Увеличение числа негативных примеров при обучении SASRec с gBCE даёт улучшение NDCG@10 порядка 9.47\% на MovieLens-1M по сравнению с лучшей конфигурацией BERT4Rec при сокращении времени обучения на ~73\%. Эксперименты на MovieLens-1M, Steam и Gowalla подтверждают, что gBCE масштабируется в широком диапазоне размеров каталога.

Для настоящей работы gSASRec имеет двойное значение. Во-первых, именно эта пара (архитектура SASRec + функция потерь gBCE) является яндексовским бейзлайном на YAMBDA~\cite{yambda2025}~--- датасете, на котором проводится эксперимент, и значит представляет собой естественную точку отсчёта при воспроизведении. Во-вторых, ключевой эмпирический результат главы~\ref{ch:experiment} (раздел~\ref{subsec:ch3:setup-gap}) состоит в том, что на YAMBDA gBCE \emph{не даёт} устойчивого улучшения над обычной (plain) BCE, и финальная конфигурация персонального скорера использует именно plain BCE. Этот результат не противоречит~\cite{gsasrec2023}, а уточняет границы применимости: эффект gBCE сильнее всего проявляется на датасетах с меньшим каталогом и меньшей долей повторных взаимодействий, в то время как музыкальный домен (см.~раздел~\ref{sec:ch1:music-domain}) обладает противоположными свойствами. Соответственно, в качестве методологически чистого бейзлайна в эксперименте используется plain SASRec~+~BCE, а отказ от gBCE подаётся не как априорное проектное решение, а как эмпирически обоснованный шаг лестницы воспроизведения яндексовского результата.

\section{Специфика музыкального домена}
\label{sec:ch1:music-domain}

\subsection{Повторное потребление и Personalized Popularity Awareness}
\label{subsec:ch1:repeated}

Музыкальный домен принципиально отличается от рекомендаций фильмов, товаров или новостей одной устойчивой особенностью: пользователи регулярно переслушивают одни и те же треки. Abbattista и соавторы~\cite{abbattista2024sequential} детально исследуют этот феномен на данных Yandex Music Event 2019 (предшественник YAMBDA, существенно меньший по объёму) и показывают, что простой бейзлайн Personalized Most Popular~--- рекомендация наиболее часто прослушиваемых данным пользователем треков~--- оказывается неожиданно сильным конкурентом нейросетевых моделей.

Авторы предлагают механизм Personalized Popularity Awareness (PPA), инъецирующий персонализированные оценки популярности в трансформерные модели семейства SASRec/BERT4Rec/gSASRec. Оценка популярности для пары пользователь--трек вычисляется как нормализованная частота прослушиваний:
\begin{equation*}
    \mathrm{pop}(u, i) = \frac{c(u, i)}{\max_{j \in \mathcal{I}_u} c(u, j)},
\end{equation*}
где $c(u, i)$~--- число прослушиваний трека~$i$ пользователем~$u$, а $\mathcal{I}_u$~--- множество прослушанных им треков. Эксперименты~\cite{abbattista2024sequential} демонстрируют улучшение метрик ранжирования на 25--70\% для всех трёх трансформерных архитектур, причём наибольший эффект наблюдается именно на музыкальных данных. Результат указывает на то, что любая последовательная рекомендательная модель в музыке должна явно или неявно учитывать паттерн повторного потребления, иначе значительная часть сигнала остаётся не использованной.

Для настоящей работы феномен повторного потребления имеет ещё одно, методологическое, измерение. В большинстве стандартных протоколов оценки последовательных рекомендеров история пользователя \emph{маскируется} при ранжировании на тестовом срезе~--- это унаследовано из протоколов оценки в доменах с малой долей повторов (фильмы, e-commerce), где переслушать <<уже потреблённый>> элемент содержательно бессмысленно. В музыке такое маскирование систематически занижает качество, поскольку значительная доля релевантных для пользователя кандидатов как раз является уже прослушанными треками. В разделе~\ref{subsec:ch3:setup-gap} показано, что снятие маски истории при ранжировании на YAMBDA даёт прирост NDCG@10 примерно в 3.2 раза по сравнению с маскированным протоколом, и именно этот шаг лестницы является решающим для закрытия setup-gap'а с яндексовским бейзлайном. Таким образом, наблюдения~\cite{abbattista2024sequential} мотивируют не только PPA как архитектурное решение, но и пересмотр eval-протокола для музыкального домена в целом.

\subsection{Роль порядка взаимодействий в музыке}
\label{subsec:ch1:order}

Дополнительное доменное наблюдение представлено в работе Greer и соавторов~\cite{greer2025beyond} на данных Amazon Music. Авторы исследуют роль архитектурных компонентов трансформера, отвечающих за порядок последовательности, и фиксируют контринтуитивный результат: удаление каузальной маски и позиционных эмбеддингов из gSASRec приводит к улучшению F1 на 28\%. Иными словами, в музыкальном домене строгий хронологический порядок взаимодействий несёт меньше полезной информации, чем в других последовательных задачах, и модель выигрывает от возможности свободно агрегировать сигнал по всей истории прослушиваний.

Этот результат содержательно перекликается с наблюдением о маскировании истории, обсуждавшимся выше: и там, и здесь стандартные индуктивные смещения трансформера последовательных рекомендаций (каузальность, позиционная информация, исключение прослушанных треков) оказываются менее эффективными в музыке, чем в типичных бенчмарках. В настоящей работе мы оставляем каузальную маску и позиционные эмбеддинги в архитектуре персонального скорера (раздел~\ref{subsec:ch3:sasrec-arch})~--- это решение продиктовано целью воспроизвести яндексовский бейзлайн на тех же гиперпараметрах,~--- однако содержательно фиксируем, что снятие маски истории при \emph{ранжировании} (но не при обучении) является самостоятельным методологическим вкладом главы~\ref{ch:experiment}, согласующимся с линией наблюдений~\cite{abbattista2024sequential, greer2025beyond}.

\section{Контентные признаки и аудиоэмбеддинги}
\label{sec:ch1:content}

\subsection{Глубокие контентные представления музыки}
\label{subsec:ch1:vandenoord}

Музыкальный домен обладает уникальным преимуществом перед другими задачами рекомендаций: аудиосигнал трека несёт богатую информацию о его характеристиках~--- жанре, темпе, настроении, инструментовке,~--- что позволяет строить содержательные представления треков даже при полном отсутствии коллаборативного сигнала. Van den Oord и соавторы~\cite{vandeoord2013deep_music} в одной из первых работ на стыке глубокого обучения и музыкальных рекомендаций предложили использовать свёрточные нейронные сети для извлечения латентных представлений из спектрограмм аудиозаписей. Ключевой результат работы состоит в том, что обученные на аудио представления коррелируют с латентными факторами коллаборативной фильтрации, а основанная на них модель способна предсказывать предпочтения пользователей даже для треков, отсутствующих в матрице взаимодействий~--- то есть служить решением задачи холодного старта по элементу.

\subsection{Аудиоэмбеддинги YAMBDA как индивидуальный сигнал}
\label{subsec:ch1:yambda-audio}

Дальнейшее развитие подхода привело к тому, что в современных датасетах музыкальных рекомендаций аудиоэмбеддинги нередко предоставляются как готовый предвычисленный признак. В частности, датасет YAMBDA~\cite{yambda2025} распространяется вместе с аудиоэмбеддингами всех треков размерности $d_a = 128$, полученными собственной аудиомоделью Яндекс.Музыки; технические детали извлечения и использования subset'а аудиоэмбеддингов в эксперименте описаны в подразделе~\ref{subsec:ch3:audio}.

С точки зрения архитектуры персонального скорера наличие готовых аудиоэмбеддингов открывает несколько ролей: они могут служить признаком для ранжирующей модели в многоступенчатой воронке, инициализацией элементных эмбеддингов трансформерных моделей или основой для рекомендации <<холодных>> треков. В контексте настоящей работы аудиоэмбеддинги выполняют существенно более узкую роль: они \emph{не} используются для обучения персонального скорера (он построен исключительно на коллаборативном сигнале~--- ID-эмбеддингах треков), но образуют пользовательский аудиопрофиль
\begin{equation}
    \bar{\mathbf{a}}_u = \frac{1}{|\mathcal{S}_u|} \sum_{i \in \mathcal{S}_u} \mathbf{a}_i,
    \label{eq:user-audio-profile}
\end{equation}
который является входом для группового агрегатора, рассматриваемого в главе~\ref{ch:groups}. Содержательно это означает, что в данной работе аудио рассматривается как уже доступный, по построению информативный индивидуальный признак, а основная исследовательская гипотеза формулируется на уровне группы: даёт ли замена ID-эмбеддингов на аудио в attention-механизме группового агрегатора устойчивое улучшение качества групповых рекомендаций. Подробное обсуждение этой гипотезы и её эмпирическая проверка проводятся в главе~\ref{ch:groups} и в разделе~\ref{sec:ch3:results}.

\section{Выводы по главе}
\label{sec:ch1:conclusion}

Проведённый обзор очерчивает три методологические опоры, на которых строится экспериментальная часть работы.

Во-первых, в качестве персонального последовательного скорера выбрана архитектура SASRec~\cite{sasrec2018} с обычной (plain) бинарной кросс-энтропией. Обоснование двусоставное. С одной стороны, замечание~\cite{petrov2022bert4rec_replicability} о воспроизводимости BERT4Rec и архитектурно нейтральный вывод~\cite{gsasrec2023} о том, что ключевой фактор качества~--- функция потерь и стратегия негативного сэмплирования, а не блок внимания, делают выбор SASRec предпочтительным с точки зрения прозрачности эксперимента. С другой стороны, проверка применимости gBCE на YAMBDA, проводимая в~разделе~\ref{subsec:ch3:setup-gap}, показывает, что на данном датасете gBCE не даёт устойчивого улучшения над plain BCE; финальная конфигурация скорера воспроизводит яндексовский бейзлайн~\cite{yambda2025} в пределах статистического шума на NDCG@10 при существенно более простой функции потерь.

Во-вторых, обзор музыкального домена~--- феномен повторного потребления и роль порядка взаимодействий~\cite{abbattista2024sequential, greer2025beyond}~--- мотивирует пересмотр стандартного eval-протокола для последовательных рекомендеров. В разделе~\ref{subsec:ch3:setup-gap} показано, что снятие маски истории при ранжировании на тестовом срезе даёт многократный прирост NDCG@10 и является решающим шагом для согласования наших чисел с яндексовским бейзлайном. Содержательно это самостоятельный методологический вклад работы: protocol-level адаптация под доменную специфику оказывается важнее, чем тонкая настройка функции потерь.

В-третьих, аудиоэмбеддинги YAMBDA~\cite{vandeoord2013deep_music, yambda2025} рассматриваются как готовый индивидуальный контентный сигнал. Их роль на индивидуальном уровне исчерпывающе исследована в литературе и в данной работе принимается как факт: персональный скорер построен на ID-эмбеддингах треков и аудио в нём не использует. Однако открытым остаётся вопрос о том, как использовать индивидуальные аудиопрофили в~\emph{групповом} агрегаторе~--- и именно этому вопросу посвящена глава~\ref{ch:groups}: классические стратегии и нейросетевые ID-based агрегаторы (AGREE, GroupIM) рассматриваются как бейзлайны, а в качестве предлагаемой архитектурной модификации вводится audio-aware attention, в котором ID-эмбеддинги пользователей заменены на их аудиопрофили~$\bar{\mathbf{a}}_u$ (формула~\eqref{eq:user-audio-profile}). Эмпирическая проверка соответствующей гипотезы проводится в главе~\ref{ch:experiment}.
